\documentclass{article}

\usepackage{
	arxiv,
	amsmath,
	amssymb,
	bm,
	cleveref
}

\title{Computing the Metric Dimension of Graphs using Satisfiability Modulo Theories}

\author{
	Fikri Yudhistira Ajiwijaya \\
	Institut Teknologi Sepuluh Nopember
	\And
	Muhammad Syifa'ul Mufid \\
	Institut Teknologi Sepuluh Nopember
}

\begin{document}
	\maketitle

	\begin{abstract}
		...\textit{blah blah blah}...
	\end{abstract}

	\section{Introduction}

	The theoretical part of this paper is divided into three major parts. The first of which consists of the heavy lifting of formally defining notations for slicing, dicing, and filtering of matrices, tuples, and sets. While it may seem trivial and verbose at first glance, it is of utmost importance due to the lack of standard notations for these basic operations. Part two and part three extensively relies on using these carefully defined operations, without which would be incomprehensible to read. This rigor is both necessary and unnegotiable. Part two consists of a quick recall on the metric dimension problem definition, then followed by the first major contribution of this paper; that is, the journey in perspective shift of metric dimension reformulated in terms of slicing, dicing, and filtering. After all these gruelling groundwork, we are finally able to introduce our great result in part three, where we formulate constraints for metric dimension in satisfiability modulo theories, through careful solution space analysis.

	...\textit{blah blah blah}...

	\section{...\textit{blah blah blah}...}

	...\textit{blah blah blah}...

	\section{Experiments and Results}

	...\textit{blah blah blah}...

	\section{Conclusions and Future Works}

	...\textit{blah blah blah}...
\end{document}
